\questao{92}{corD}{Energia e calor}{Potência térmica e conversão de unidades}{Com conta}
\begin{enunciado}
% gerado por ferramentas/md2tex_enunciado.py a partir de questoes/q092.md (não editar à mão)
O quadro lista o resultado da simulação feita por uma engenheira para dimensionar as fontes de calor no interior de um escritório onde trabalham 20 pessoas. Todos os equipamentos e as pessoas dissipam calor de acordo com os valores no quadro. Para manter constante a temperatura do escritório, ela precisa de um aparelho de ar-condicionado o mais econômico possível. Considere as paredes isolantes e o fator de conversão 1\,Btu/h = 0,3\,W.\par
\begin{center}\renewcommand{\arraystretch}{1.15}
\begin{tabular}{ccc}
\toprule
\textbf{Fonte de calor} & \textbf{Quantidade} & \textbf{Potência individual (W)}\\
\midrule
Lâmpada & 10 & 100\\
Computador & 20 & 150\\
Geladeira & 2 & 500\\
Pessoa adulta & 20 & 200\\
\bottomrule
\end{tabular}\end{center}
Qual é a capacidade de refrigeração, em British thermal unit por hora, do aparelho selecionado?\par
\begin{altlinha}
\alth{A}{7\,000\,Btu/h}%
\alth{B}{10\,000\,Btu/h}%
\alth{C}{12\,000\,Btu/h}%
\alth{D}{24\,000\,Btu/h}%
\alth{E}{30\,000\,Btu/h}
\end{altlinha}

\end{enunciado}

\quer{qual deve ser a capacidade do ar-condicionado, em Btu/h, para tirar todo o calor produzido dentro do escritório.}

\begin{ideia}
\textbf{Potência} é a quantidade de energia por segundo. O joule (J) é a unidade de energia, e 1 watt (W) é 1 joule por segundo.
Uma lâmpada de 100\un{W} joga 100 joules de calor na sala a cada segundo.
Uma pessoa também aquece a sala: nosso corpo libera calor o tempo todo.

Para a temperatura \textbf{não subir nem descer}, o ar-condicionado precisa tirar calor \textbf{no mesmo ritmo} em que o calor é produzido.
O enunciado diz que as paredes são \textbf{isolantes}: não entra nem sai calor por elas. Então, todo o calor que o aparelho
precisa tirar é o que é produzido lá dentro.

\textbf{Pense em uma pia com a torneira aberta.} Para o nível da água ficar parado, o ralo precisa escoar exatamente o que a
torneira despeja. As lâmpadas, os computadores, as geladeiras e as pessoas são as ``torneiras de calor''. O ar-condicionado é o ``ralo''.
\end{ideia}

\passoapasso
\npasso{Calcule o calor produzido por cada grupo.}
Multiplique a quantidade pela potência de cada um:
\begin{conta}
\renewcommand{\arraystretch}{1.25}
\begin{tabular}{lcl}
Lâmpadas: & $10 \times 100 =$ & $1\,000\un{W}$\\
Computadores: & $20 \times 150 =$ & $3\,000\un{W}$\\
Geladeiras: & $2 \times 500 =$ & $1\,000\un{W}$\\
Pessoas: & $20 \times 200 =$ & $4\,000\un{W}$
\end{tabular}
\end{conta}

\npasso{Some tudo.}
\begin{conta}
$1\,000 + 3\,000 + 1\,000 + 4\,000 = 9\,000\un{W}$
\end{conta}
O escritório produz 9\,000\un{W} de calor. O aparelho precisa retirar esses 9\,000\un{W}.

\npasso{Troque a unidade: de watt para Btu/h.}
O enunciado informa que 1\un{Btu/h} vale 0,3\un{W}. Queremos saber \textbf{quantos ``pedaços'' de 0,3\un{W} cabem em 9\,000\un{W}}.
Isso é uma divisão:
\begin{conta}
$9\,000 \div 0{,}3$
\end{conta}
Dividir por número com vírgula atrapalha. Truque: multiplique os dois números por 10 (assim 0,3 vira 3 e a vírgula some). O resultado da divisão não muda:
\begin{conta}
$9\,000 \div 0{,}3 \;=\; 90\,000 \div 3 \;=\; 30\,000\un{Btu/h}$
\end{conta}
\confere{$30\,000 \times 0{,}3 = 9\,000\un{W}$. Bateu.}

\npasso{Entenda o ``mais econômico possível''.}
O aparelho mais econômico é o \textbf{menor} que ainda dá conta do serviço. Um aparelho de 24\,000\un{Btu/h} retira só
$24\,000 \times 0{,}3 = 7\,200\un{W}$. É menos que 9\,000\un{W}, e a sala esquentaria. Então o aparelho certo é o de 30\,000\un{Btu/h}.
(Para fazer $24\,000 \times 0{,}3$: $24\,000 \times 3 = 72\,000$; como 0,3 tem uma casa depois da vírgula, 72\,000 vira 7\,200.)

\resposta{E}{O escritório produz 9\,000\un{W} de calor, que equivalem a 30\,000\un{Btu/h}.}

\begin{outras}
\altl{A}, \altl{B}, \altl{C} e \alt{D} Todos esses aparelhos retiram menos de 9\,000\un{W}. O maior deles (24\,000\un{Btu/h}) retira
7\,200\un{W}. Sobraria calor na sala e a temperatura subiria.
\end{outras}

\begin{pegadinha}[Onde muita gente erra]
\begin{itemize}
  \item \textbf{Esquecer as pessoas.} Elas são a maior fonte de calor do quadro (4\,000\un{W}).
  \item \textbf{Esquecer de multiplicar pela quantidade.} São 10 lâmpadas, não uma.
  \item \textbf{Multiplicar por 0,3 em vez de dividir.} Dica para não errar: 1\un{Btu/h} é ``menor'' que 1\un{W}. Então o número em
        Btu/h tem de ser \textbf{maior} que o número em watts.
\end{itemize}
\end{pegadinha}
